\centerline{\Huge Условия и решения задач.}
\centerline{\bf \Huge 6 класс}
\problem{6.1} На доске написаны две суммы:\par
1 + 22 + 333 + 4444 + 55555 + 666666 +7777777 + 88888888 + 999999999 \par
9 + 98 + 987 + 9876 + 98765 + 987654 + 9876543 + 98765432 + 987654321 \\ 
Определите, какая из них больше (или они равны). \par

\textbf{Решение.} 
Запишем обе суммы столбиком, причём вторую для большей наглядности в обратном порядке. В обеих суммах в разряде единиц будут складываться цифры от 1 до 9, в разряде десятков — цифры от 2 до 9, в разряде сотен — цифры от 3 до 9 и так далее. И цифра, получающаяся в каждом разряде, и перенос в старший разряд окажутся в обоих случаях одинаковыми. Поэтому и результат сложения будет один и тот же.
\begin{figure}[h] 
    \includegraphics[scale=0.3]{Региональная Олимпиада в Калмыкии 2023/Решения/pic1.png}
\end{figure} \\
\\

\textbf{Критерии.} \\
\textit{7 баллов} $-$ Полностью верное решение. \\
\textit{0-2 балла} $-$ Есть какие-то продвижения, которые могут привести к решению. \\ 
\textit{0 баллов} $-$ Решение отсутствует. Решение неверное, продвижения отсутствуют. \\ 

\problem{6.2} На какую цифру оканчивается число $2023^{2022}$? \par

\textbf{Первое решение.}  Заметим, что достаточно лишь следить за тем, как изменяется последняя цифра числа $3^{2022}$, а не $2023^{2022}$, так как последняя цифра степени числа изменяется только лишь благодаря последней цифре самого числа. На языке арифметики остатков это означает, что для нахождения последней цифры числа, нужно \newpage рассмотреть его остаток при делении на 10. Посмотрим на то, какие последние цифры бывают у степеней тройки: $3^1$ оканчивается на 3, $3^2$ оканчивается на 9, $3^3$ оканчивается на 7, $3^4$ оканчивается на 1, $3^5$ оканчивается снова на 3. Видим, что мы зациклились, то есть последние цифры степеней тройки повторяются через 4. Значит, что у $3^{2022}$ такая же последняя цифра как и у $3^2$. То есть 9. \par 

\textbf{Второе решение.} Найти последнюю цифру это то же самое, что найти остаток при делении на 10. По теореме Эйлера имеем, что $2023^{\varphi (10)} \equiv 3^{\varphi (10)} = 3^4 \equiv 1 \ (mod \ 10)$. То есть $2023^{2022} \ \equiv \ 3^{2022} \  \equiv \  3^{2020}3^2 \ \equiv \ (3^4)^{505}3^2 \ \equiv \ 3^2 = 9 \ (mod \ 10)$ \\ 

\textbf{Критерии.} \\
\textit{7 баллов} $-$ Полностью верное решение. \\
\textit{4 балла} $-$ Доказано, что последняя цифра степени тройки повторяется каждые 4 раза. \\ 
\textit{2 балла} $-$ Доказано, что вместо $2023^{2022}$ можно смотреть на $3^{2022}$.  \\
\textit{0-2 балла} $-$ Есть какие-то продвижения, которые могут привести к решению. \\ 
\textit{0 баллов} $-$ Решение отсутствует. Решение неверное, продвижения отсутствуют. \\ 

\problem{6.3} В Древней Спарте выживали сильнейшие. Поэтому 300 спартанцев решили выяснить, кто из них достойный воин. Они разбились по парам и сделали тренировочный бой. Потом опять разбились по парам (возможно, уже по-другому) и снова провели тренировочный бой. Тот, кто выигрывал оба боя, считался достойным воином, а остальных сбрасывали со скалы. Какое минимальное количество воинов могло остаться в Спарте? \par
Известно, что все 300 спартанцев были разной силы, а выигрывал в бою всегда тот, у кого силы больше.\par

\textbf{Решение.} Самый сильный спартанец обязательно выживет.
  Покажем, что может быть ровно один выживший спартанец. Пронумеруем спартанцев по возрастанию силы от 1 до 300. В первый раз проведём поединки так: \\ 1 $-$ 2, \\ 3 $-$ 4, \\ 5 $-$ 6, \\ ..., \\ 299 $-$ 300. \\ А во второй раз так: \\ 300 $-$ 1, \\ 2 $-$ 3, \\ 4 $-$ 5, \\ ..., \\ 298 $-$ 299. \\ Тогда каждый, кроме самого сильного, в одном из туров проиграет. \\ 

\textbf{Критерии.} \\
\textit{7 баллов} $-$ Полностью верное решение. \\
\textit{2 балла} $-$ Доказано, что самый сильный спартанец всегда выживет. \\
\textit{0-2 балла} $-$ Есть какие-то продвижения, которые могут привести к решению. \\ 
\textit{0 баллов} $-$ Решение отсутствует. Решение неверное, продвижения отсутствуют. \\ 

%https://problems.ru/view_problem_details_new.php?id=64440

\problem{6.4} Лера, Аркадий и Владимир записали на доске по одному числу $-$ 0, 2 и 5 соответственно. Спустя минуту каждый из них одновременно стирает свое число и записывает вместо него сумму двух других(то есть Лера запишет 7, Аркадий 5, а Владимир 2). Спустя ещё минуту они повторяют это действие и так несколько раз. Могла ли сумма чисел Леры, Аркадия и Владимира в какой-то момент стать равной 20232022? \par 

\textbf{Решение.} Давайте будем следить за суммой всех трех чисел. Если изначально были числа $a, \ b$ и $c$, то через минуту вместо них будут $b+c, \ a+c,$ и $a+b$ соответсвенно. В первый раз сумма трех чисел была равна $a+b+c$, а через минуту уже $(b+c) + (a+c) + (a+b) = 2(a+b+c)$. То есть каждую минуту сумма всех трёх чисел удваивается. Тогда через $n$ минут сумма будет равна $2^n(0+2+5) = 7\cdot2^n$. Но такое число не может быть равно $20232022$ по разным причинам. Например, 20232022 не делится на 7 и на 4.\\

\textbf{Критерии.} \\
\textit{7 баллов} $-$ Полностью верное решение. \\
\textit{5 баллов} $-$ Доказано, что сумма трех чисел увеличивается в два раза с каждым разом. \\
\textit{0-2 балла} $-$ Есть какие-то продвижения, которые могут привести к решению. \\ 
\textit{0 баллов} $-$ Решение отсутствует. Решение неверное, продвижения отсутствуют. \\ 

\problem{6.5} Очир выписал на доску числа $1, \dfrac{1}{2}, \dfrac{1}{3}, \dfrac{1}{4}, \dotsc, \dfrac{1}{2023}$. За один ход он может стереть любые 
\\
\\
два числа $a$ и $b$ и заменить их на число $ab+a+b$. Какое число останется, если Очир проделает 2022 такие операции? \par 

\textbf{Решение.} Пусть в некоторый момент на доске написаны числа $a, b, ... , z.$ Рассмотрим произведение $ P = (a + 1)(b + 1)...(z + 1).$  После замены пары чисел $a, b$ на  $ab + a + b $ в указанном произведении  $(a + 1)(b + 1)$ заменится на  $ab + a + b + 1 = (a + 1)(b + 1)$.  Таким образом, операция величину $P$ не меняет. Отсюда следует, что число, которое остаётся в конце, не зависит от порядка выполнения операций и равно: \\

$\Biggl(1+1\Biggr)\Biggl(\dfrac{1}{2} + 1\Biggr)\Biggl(\dfrac{1}{3} + 1\Biggr) \dotsc \Biggl(\dfrac{1}{2023} + 1\Biggr) = \dfrac{2}{1} \cdot \dfrac{3}{2} \cdot \dfrac{4}{3} \dotsc \dfrac{2023}{2022} = 2023$ \\ 

\textbf{Критерии.} \\
\textit{7 баллов} $-$ Полностью верное решение. \\
\textit{5 баллов} $-$ Доказано, что сохраняется произведение чисел увеличенных на 1. \\
\textit{0-2 балла} $-$ Есть какие-то продвижения, которые могут привести к решению. \\ 
\textit{0 баллов} $-$ Решение отсутствует. Решение неверное, продвижения отсутствуют. \\ 

%https://problems.ru/view_problem_details_new.php?id=35225